% GENERATED FILE. Edit canonical lessons, then run npm run generate:books.
% canonical-source-hash: fnv1a64:b1fa76b5ad404554
% canonical-lessons: PT-C306-evacuar, PT-C306-resgatar, PT-C306-resgate, PT-C306-sobreviver, PT-C306-sobrevivente

\chapter{Todos salvos}
\label{ch:pt-todos-salvos}
\hlchaptermodality{306}

\begin{chapteropening}
\textbf{By the end of this chapter:} \emph{I can tell how people were evacuated and rescued, and who survived.}

Complete the last lesson of chapter 306: I can tell how people were evacuated and rescued, and who survived.
\end{chapteropening}

\section[evacuar]{evacuar — to evacuate}

\label{lesson:PT-C306-evacuar}



\begin{quote}
Before the new one: say the Portuguese for the first-aider, then the Portuguese for the stretcher.
\end{quote}

\subsection*{You'll want to know: evacuar}
\textbf{evacuar} — \textquotedblleft{}to evacuate\textquotedblright{}.

\begin{quote}
Evacuaram a escola por causa do fumo.
\end{quote}

\begin{quote}
\emph{They evacuated the school because of the smoke.}
\end{quote}

Takes a direct object: \emph{evacuar o prédio}.

\begin{cousinweb}
\emph{evacuar} is a learned borrowing from Latin \emph{ēvacuāre} \textquotedblleft{}to empty out\textquotedblright{} (EtymDB), built on \emph{vacuus} \textquotedblleft{}empty\textquotedblright{}.
\end{cousinweb}

\begin{grammarlens}[title={What you've built}]
One more word for telling, explaining and describing at B1.
\end{grammarlens}

\subsection*{Your turn}
\begin{itemize}
\raggedright
  \item \emph{Say it:} \emph{evacuar}
  \item \emph{Say it:} \emph{evacuar}, once more
  \item \emph{Say it:} a sentence of your own with \emph{evacuar}
  \item \emph{Recall:} say the Portuguese for the first-aider, then the Portuguese for the stretcher, then say \emph{evacuar} again
\end{itemize}

\begin{tcolorbox}[breakable,colback=teal!4,colframe=teal!35!black,title={Before you move on}]
What does evacuar mean? (\textquotedblleft{}To evacuate\textquotedblright{}.) Say it once more.
\end{tcolorbox}
\section[resgatar]{resgatar — to rescue; to ransom}

\label{lesson:PT-C306-resgatar}



\begin{quote}
Before the new one: say the Portuguese for the stretcher, then the Portuguese for to evacuate.
\end{quote}

\subsection*{You'll want to know: resgatar}
\textbf{resgatar} — \textquotedblleft{}to rescue; to ransom\textquotedblright{}.

\begin{quote}
Os bombeiros resgataram um gato da árvore.
\end{quote}

\begin{quote}
\emph{The firefighters rescued a cat from the tree.}
\end{quote}

Also to cash in: \emph{resgatar um prémio}.

\begin{cousinweb}
\emph{resgatar} is built on \emph{resgate} (EtymDB), from Spanish \emph{rescatar} \textquotedblleft{}to rescue\textquotedblright{}, which goes back to Vulgar Latin \emph{recaptāre}, from Latin \emph{captāre} \textquotedblleft{}to catch\textquotedblright{}.
\end{cousinweb}

\begin{grammarlens}[title={What you've built}]
One more word for telling, explaining and describing at B1.
\end{grammarlens}

\subsection*{Your turn}
\begin{itemize}
\raggedright
  \item \emph{Say it:} \emph{resgatar}
  \item \emph{Say it:} \emph{resgatar}, once more
  \item \emph{Say it:} a sentence of your own with \emph{resgatar}
  \item \emph{Recall:} say the Portuguese for the stretcher, then the Portuguese for to evacuate, then say \emph{resgatar} again
\end{itemize}

\begin{tcolorbox}[breakable,colback=teal!4,colframe=teal!35!black,title={Before you move on}]
What does resgatar mean? (\textquotedblleft{}To rescue; to ransom\textquotedblright{}.) Say it once more.
\end{tcolorbox}
\section[o resgate]{o resgate — the rescue; the ransom}

\label{lesson:PT-C306-resgate}



\begin{quote}
Before the new one: say the Portuguese for to evacuate, then the Portuguese for to rescue.
\end{quote}

\subsection*{You'll want to know: o resgate}
\textbf{o resgate} — \textquotedblleft{}the rescue; the ransom\textquotedblright{}.

\begin{quote}
O resgate durou toda a noite.
\end{quote}

\begin{quote}
\emph{The rescue went on all night.}
\end{quote}

In a kidnapping, the money: \emph{pedir um resgate}.

\begin{cousinweb}
\emph{resgate} was borrowed from Spanish \emph{rescate} (EtymDB), from \emph{rescatar}, which goes back to Vulgar Latin \emph{recaptāre}, from Latin \emph{captāre} \textquotedblleft{}to catch\textquotedblright{}.
\end{cousinweb}

\begin{grammarlens}[title={What you've built}]
One more word for telling, explaining and describing at B1.
\end{grammarlens}

\subsection*{Your turn}
\begin{itemize}
\raggedright
  \item \emph{Say it:} \emph{o resgate}
  \item \emph{Say it:} \emph{o resgate}, once more
  \item \emph{Say it:} a sentence of your own with \emph{o resgate}
  \item \emph{Recall:} say the Portuguese for to evacuate, then the Portuguese for to rescue, then say \emph{o resgate} again
\end{itemize}

\begin{tcolorbox}[breakable,colback=teal!4,colframe=teal!35!black,title={Before you move on}]
What does o resgate mean? (\textquotedblleft{}The rescue; the ransom\textquotedblright{}.) Say it once more.
\end{tcolorbox}
\section[sobreviver]{sobreviver — to survive}

\label{lesson:PT-C306-sobreviver}



\begin{quote}
Before the new one: say the Portuguese for to rescue, then the Portuguese for the rescue.
\end{quote}

\subsection*{You'll want to know: sobreviver}
\textbf{sobreviver} — \textquotedblleft{}to survive\textquotedblright{}.

\begin{quote}
Todos sobreviveram ao acidente.
\end{quote}

\begin{quote}
\emph{Everyone survived the accident.}
\end{quote}

\emph{Sobreviver a} + what you lived through.

\begin{cousinweb}
\emph{sobreviver} comes from Latin \emph{supervīvere} \textquotedblleft{}to outlive, to survive\textquotedblright{} (EtymDB), from \emph{super} \textquotedblleft{}over\textquotedblright{} and \emph{vīvere} \textquotedblleft{}to live\textquotedblright{}.
\end{cousinweb}

\begin{grammarlens}[title={What you've built}]
One more word for telling, explaining and describing at B1.
\end{grammarlens}

\subsection*{Your turn}
\begin{itemize}
\raggedright
  \item \emph{Say it:} \emph{sobreviver}
  \item \emph{Say it:} \emph{sobreviver}, once more
  \item \emph{Say it:} a sentence of your own with \emph{sobreviver}
  \item \emph{Recall:} say the Portuguese for to rescue, then the Portuguese for the rescue, then say \emph{sobreviver} again
\end{itemize}

\begin{tcolorbox}[breakable,colback=teal!4,colframe=teal!35!black,title={Before you move on}]
What does sobreviver mean? (\textquotedblleft{}To survive\textquotedblright{}.) Say it once more.
\end{tcolorbox}
\section[o sobrevivente]{o sobrevivente — the survivor}

\label{lesson:PT-C306-sobrevivente}



\begin{quote}
Before the new one: say the Portuguese for the rescue, then the Portuguese for to survive.
\end{quote}

\subsection*{You'll want to know: o sobrevivente}
\textbf{o sobrevivente} — \textquotedblleft{}the survivor\textquotedblright{}.

\begin{quote}
O único sobrevivente contou tudo à polícia.
\end{quote}

\begin{quote}
\emph{The only survivor told the police everything.}
\end{quote}

The same form for a woman: \emph{a sobrevivente}.

\begin{cousinweb}
\emph{sobrevivente} is built on \emph{sobreviver} \textquotedblleft{}to survive\textquotedblright{}, from Latin \emph{supervīvere} (EtymDB), with the ending \emph{-ente}.
\end{cousinweb}

\begin{grammarlens}[title={What you've built}]
One more word for telling, explaining and describing at B1.
\end{grammarlens}

\subsection*{Your turn}
\begin{itemize}
\raggedright
  \item \emph{Say it:} \emph{o sobrevivente}
  \item \emph{Say it:} \emph{o sobrevivente}, once more
  \item \emph{Say it:} a sentence of your own with \emph{o sobrevivente}
  \item \emph{Recall:} say the Portuguese for the rescue, then the Portuguese for to survive, then say \emph{o sobrevivente} again
\end{itemize}

\begin{tcolorbox}[breakable,colback=teal!4,colframe=teal!35!black,title={Before you move on}]
What does o sobrevivente mean? (\textquotedblleft{}The survivor\textquotedblright{}.) Say it once more.
\end{tcolorbox}
